\section*{第 \ref{ch:g10:universal-gravitation} 章——万有引力与重力}

\begin{solution}{exo:g10:universal-gravitation:1}
\emph{1.} $F = \num{6.67e-11} \times \dfrac{60 \times 60}{0.80^2}
= \dfrac{\num{2.40e-7}}{0.64} \approx \qty{3.8e-7}{N}$.

\emph{2.} 蚊子的\omterm{def:g10:universal-gravitation:weight}{重力}为
$P = \num{2.5e-6} \times 9.81 \approx \qty{2.5e-5}{N}$——约比两名舞者之间的
相互吸引力\emph{大} $65$ 倍。人与人之间的
\omterm{def:g10:universal-gravitation:force}{万有引力}比蚊子的\omterm{def:g10:universal-gravitation:weight}{重力}还小。
\end{solution}

\begin{solution}{exo:g10:universal-gravitation:2}
\emph{1.} $P = mg = 55 \times 9.81 \approx \qty{5.4e2}{N}$.

\emph{2.} 她的质量不变，仍为 \qty{55}{kg}（所含
物质的多少不变）。她的\omterm{def:g10:universal-gravitation:weight}{重力}变为
$P = 55 \times 1.62 \approx \qty{89}{N}$——为原来的六分之一。
\end{solution}

\begin{solution}{exo:g10:universal-gravitation:3}
\emph{1.} $g_{\text{火星}} = \dfrac{G M}{R^2}
= \dfrac{\num{6.67e-11} \times \num{6.42e23}}{(\num{3.39e6})^2}
= \dfrac{\num{4.28e13}}{\num{1.15e13}} \approx \qty{3.73}{N/kg}$.

\emph{2.} 在火星上：$P = 900 \times 3.73 \approx \qty{3.4e3}{N}$；在
地球上：$P = 900 \times 9.81 \approx \qty{8.8e3}{N}$——探测车
在火星上的\omterm{def:g10:universal-gravitation:weight}{重力}约小 $2.6$ 倍，而质量相同。
\end{solution}

\begin{solution}{exo:g10:universal-gravitation:4}
\[
F = \num{6.67e-11} \times
\frac{\num{5.97e24} \times \num{7.35e22}}{(\num{3.84e8})^2}
= \frac{\num{2.93e37}}{\num{1.47e17}}
\approx \qty{2.0e20}{N}.
\]
月球对地球施加一个与之\emph{大小相同}的力，
\qty{2.0e20}{N}，方向相反：\omterm{def:g10:universal-gravitation:force}{万有引力}是相互的。
\end{solution}

\begin{solution}{exo:g10:universal-gravitation:5}
\emph{1.} 距离加倍，$d^2$ 变为原来的 $4$ 倍，因此
$F = F_0/4$.

\emph{2.} 距离减半：$F = 4F_0$。

\emph{3.} 两物体的质量均加倍，$m_1 m_2$ 变为原来的 $4$ 倍，因此
$F = 4F_0$.

\emph{4.} 分子 $\times 3$，分母 $\times 9$：
$F = F_0/3$.
\end{solution}

\begin{solution}{exo:g10:universal-gravitation:6}
\emph{1.} $F = \num{6.67e-11} \times
\dfrac{(\num{5.0e8})^2}{100^2}
= \num{6.67e-11} \times \dfrac{\num{2.5e17}}{\num{1.0e4}}
\approx \qty{1.7e3}{N}$.

\emph{2.} $m = F/g = 1667/9.81 \approx \qty{1.7e2}{kg}$：两艘船
以相当于一辆大型摩托车\omterm{def:g10:universal-gravitation:weight}{重力}的力相互吸引。

\emph{3.} $a = F/m = \num{1.7e3}/\num{5.0e8} \approx
\qty{3.3e-6}{m/s^2}$。即使不受阻碍，油轮也需要约
5 天才能达到步行速度；实际上，水的阻力和系泊
完全掩盖了这一引力。
\end{solution}

\begin{solution}{exo:g10:universal-gravitation:7}
\emph{1.} 托盘天平在两地的示数均为 \qty{3.0}{kg}：袋子
与标准砝码的\omterm{def:g10:universal-gravitation:weight}{重力}都除以同一因子，
因此比较结果不变。弹簧测力计所示的是
\omterm{def:g10:universal-gravitation:weight}{重力}：$3.0 \times 9.81 \approx \qty{29}{N}$，在地球上
$3.0 \times 1.62 \approx \qty{4.9}{N}$，在月球上。

\emph{2.} 天平：它测量的是质量，这正是顾客
所购买的东西。若弹簧测力计在地球上以“\omterm{def:g10:orders-of-magnitude:unit}{千克}”标定，
则同一袋在月球上的示数为 $4.9/9.81 \approx \qty{0.50}{kg}$——店主
就会以一袋的价钱送出六袋。
\end{solution}

\begin{solution}{exo:g10:universal-gravitation:8}
\emph{1.} 到地心的距离：
$d = R + h = \num{6.37e6} + \num{4.0e5} = \qty{6.77e6}{m}$。于是
$g = \dfrac{\num{6.67e-11} \times \num{5.97e24}}{(\num{6.77e6})^2}
= \dfrac{\num{3.98e14}}{\num{4.58e13}} \approx \qty{8.69}{N/kg}$.

\emph{2.} $8.69/9.81 \approx 89\%$，约为地表值的。

\emph{3.} 高空的\omterm{def:g10:universal-gravitation:force}{万有引力}几乎与地面一样强：宇航员
漂浮，是因为空间站及其中的一切都绕地球\emph{一起
下落}（自由落体），而不是因为
引力已经消失。
\end{solution}

\begin{solution}{exo:g10:universal-gravitation:9}
由 $g = GM/R^2$，
\[
M = \frac{g R^2}{G}
= \frac{8.87 \times (\num{6.05e6})^2}{\num{6.67e-11}}
= \frac{\num{3.25e14}}{\num{6.67e-11}}
\approx \qty{4.87e24}{kg},
\]
约为地球质量的 $0.8$ 倍——金星
与地球在大小和质量上几乎相同。
\end{solution}

\begin{solution}{exo:g10:universal-gravitation:10}
\emph{1.} 太阳：
$F = \dfrac{\num{6.67e-11} \times \num{1.99e30}}{(\num{1.496e11})^2}
\approx \qty{5.9e-3}{N}$。月球：
$F = \dfrac{\num{6.67e-11} \times \num{7.35e22}}{(\num{3.84e8})^2}
\approx \qty{3.3e-5}{N}$.

\emph{2.} 太阳的拉力约为月球的
$\num{5.9e-3}/\num{3.3e-5} \approx 180$ 倍。

\emph{3.} 地球直径在地月距离中所占的\emph{比例}（$\approx 3\%$）远大于它在日地距离中所占的比例，
它在日地距离中所占的比例，
（$\approx 0.01\%$），因此月球引力在地球面向月球的一侧与背向月球的一侧之间的差异
要大得多。潮汐正是由这一差异引起的——
在这一点上，月球占了上风。
\end{solution}

\begin{solution}{exo:g10:universal-gravitation:11}
\emph{1.} $g' = \dfrac{G (2M)}{(2R)^2} = \dfrac{2}{4}\,\dfrac{GM}{R^2}
= \dfrac{g}{2} \approx \qty{4.9}{N/kg}$：半径加倍的效应压过了
质量加倍的效应。

\emph{2.} 需要 $\dfrac{GM}{R'^2} = 2\,\dfrac{GM}{R^2}$，于是
$R'^2 = R^2/2$，并 $R' = R/\sqrt{2} \approx \qty{4.5e6}{m}$——
约 \qty{4500}{km}。
\end{solution}

\begin{solution}{exo:g10:universal-gravitation:12}
\emph{1.} 每\omterm{def:g10:orders-of-magnitude:unit}{千克}所受引力相等：
$\dfrac{G M_{\text{地球}}}{x^2} = \dfrac{G M_{\text{月球}}}{(D-x)^2}$。
交叉相乘并除以 $G$：
$\left(\dfrac{x}{D-x}\right)^2 = \dfrac{M_{\text{地球}}}{M_{\text{月球}}}$。

\emph{2.} $\dfrac{x}{D-x}
= \sqrt{\dfrac{\num{5.97e24}}{\num{7.35e22}}} = \sqrt{81.2} \approx
9.0$，于是 $x = 9.0\,(D - x)$，即 $x = \dfrac{9.0}{10.0}\,D
\approx 0.90 \times \num{3.84e8} \approx \qty{3.5e8}{m}$。平衡点
位于由地球至月球路程的 $90\%$ 处：地球的引力
几乎在整个行程中都占主导。
\end{solution}

\begin{solution}{exo:g10:universal-gravitation:13}
\emph{1.} 展开得 $R^2 + 2Rh + h^2 = R^2 + x^2$；由于 $h$ 远小于
$R$，略去 $h^2$，得 $2Rh \approx x^2$，于是
$h \approx \dfrac{x^2}{2R}
= \dfrac{(\num{8.0e3})^2}{2 \times \num{6.37e6}}
= \dfrac{\num{6.4e7}}{\num{1.27e7}} \approx \qty{5.0}{m}$.

\emph{2.} 在一秒内抛体下落 \qty{4.9}{m}——几乎
恰好等于在 \qty{8.0}{km} 内地面弯离的 \qty{5.0}{m}。因此
当 $v \approx \qty{8.0}{km/s}$ 时，下落永远赶不上地面：
抛体便绕地球运行。

\emph{3.} $T = \dfrac{2\pi R}{v}
= \dfrac{2\pi \times \num{6.37e6}}{\num{8.0e3}}
\approx \qty{5.0e3}{s} \approx \qty{83}{min}$——接近
$\qty{90}{min}$ 左右的真实近地\omterm{def:g10:universal-gravitation:satellite}{轨道}\omterm{def:g10:universal-gravitation:satellite}{卫星}周期，它们飞得
略高一些。
\end{solution}

\begin{solution}{exo:g10:universal-gravitation:14}
\emph{1.} $g = \dfrac{\num{6.67e-11} \times \num{2.8e30}}
{(\num{1.2e4})^2} = \dfrac{\num{1.87e20}}{\num{1.44e8}}
\approx \qty{1.3e12}{N/kg}.$

\emph{2.} $P = 70 \times \num{1.3e12} \approx \qty{9.1e13}{N}$，而在地球上
仅为 $\qty{6.9e2}{N}$，约为其 $\num{1.3e11}$ 倍。
在那里，任何由原子构成的结构都无法保持完整。

\emph{3.} $P = \num{1.0e-6} \times \num{1.3e12} \approx
\qty{1.3e6}{N}$。在地球上，这相当于质量为
$m = \num{1.3e6}/9.81 \approx \qty{1.3e5}{kg}$ 的物体的重量——一粒沙子的
重量竟与一台百吨机车相当。
\end{solution}

\begin{solution}{exo:g10:universal-gravitation:15}
\emph{1.} 距离变为原来的 $60$ 倍，因此每\omterm{def:g10:orders-of-magnitude:unit}{千克}所受的引力
减弱为原来的 $60^2 = 3600$ 分之一。

\emph{2.} 同理，一秒内的下落高度按相同比例变化：
$4.9/3600 \approx \qty{1.4e-3}{m}$——约 \qty{1.4}{mm}。

\emph{3.} $T = 27.3 \times \qty{86400}{s} = \qty{2.36e6}{s}$，于是
$v = \dfrac{2\pi \times \num{3.84e8}}{\num{2.36e6}} \approx
\qty{1.02e3}{m/s}$，并得一秒内的距离 $x = \qty{1.02e3}{m}$。下落
高度为 $h \approx \dfrac{x^2}{2d}
= \dfrac{(\num{1.02e3})^2}{2 \times \num{3.84e8}}
\approx \qty{1.4e-3}{m}$.

\emph{4.} 两者相符：月球每秒朝地球下落的量，
恰与以落苹果校准的平方反比引力所预言的相同。
牛顿由此断定，使月球保持绕行的力
\emph{就是}使苹果下落的那种力——引力是
普遍存在的。
\end{solution}

\begin{solution}{pb:g10:universal-gravitation:1}
\textbf{1.} $F = \num{6.67e-11} \times
\dfrac{\num{7.3e6} \times 55}{100^2} \approx \qty{2.7e-6}{N}$.

\textbf{2.} 一粒沙子的重量为
$\num{1.0e-6} \times 9.81 \approx \qty{9.8e-6}{N}$：整个埃菲尔
铁塔对游客的引力，比一粒沙子的重量\emph{小}
约 $3.7$ 倍。

\textbf{3.} (a) 距离 $\times 10$：引力除以 $100$。
(b) 两质量均 $\times 1000$：引力乘以 $10^6$。
(c) 质量 $\times 1000$，距离也 $\times 1000$：
$10^6 / 10^6$——引力不变。

\textbf{4.} $F = \dfrac{\num{6.67e-11} \times \num{5.97e24} \times
1.00}{(\num{6.37e6})^2} \approx \qty{9.81}{N}$：一\omterm{def:g10:orders-of-magnitude:unit}{千克}
物体的重量——这恰好就是
\omterm{def:g10:universal-gravitation:fieldstrength}{重力加速度} $g = \qty{9.81}{N/kg}$。

\textbf{5.} 引力\omterm{def:g10:orders-of-magnitude:unit}{单位}质量上的作用很微弱，但它只会
\emph{吸引}，因此地球全部 $\num{5.97e24}$ \omterm{def:g10:orders-of-magnitude:unit}{千克}物质的
引力都朝同一方向叠加——天文质量胜过邻近的物体，
大出二十二个\omterm{def:g10:orders-of-magnitude:oom}{数量级}。

\textbf{6.} $F = \num{6.67e-11} \times \dfrac{158 \times 0.73}{0.23^2}
= \num{6.67e-11} \times \dfrac{115.3}{0.0529}
\approx \qty{1.5e-7}{N}$。

\textbf{7.} 小球的\omterm{def:g10:universal-gravitation:weight}{重力}为
$0.73 \times 9.81 \approx \qty{7.2}{N}$——约为待测引力的 $\num{5e7}$ 倍
。普通托盘天平无法分辨五千万分之一；
一根细扭丝在纳牛量级的力矩作用下会明显扭转，
便能测出。

\textbf{8.} $G = \dfrac{F d^2}{m_1 m_2}
= \dfrac{\num{1.47e-7} \times 0.0529}{115.3}
\approx \qty{6.74e-11}{N.m^2/kg^2}$——与现代值的偏差约在 $1\%$ 以内
（现代值为 \num{6.67e-11}），该结果来自 1798 年的一间木棚。

\textbf{9.} 由 $g = G M / R^2$：
\[
M = \frac{g R^2}{G}
= \frac{9.81 \times (\num{6.37e6})^2}{\num{6.67e-11}}
\approx \qty{5.97e24}{kg}.
\]
六亿亿亿\omterm{def:g10:orders-of-magnitude:unit}{千克}：地球被称出来了。

\textbf{10.} $V = \frac43 \pi R^3 = \frac43 \pi
(\num{6.37e6})^3 \approx \qty{1.08e21}{m^3}$，因此
$\rho = \num{5.97e24}/\num{1.08e21} \approx \qty{5.5e3}{kg/m^3}$ ---
是地表岩石密度的两倍。地球内部必定藏着
比花岗岩致密得多的物质：地球有一个很重的（铁质）地核。

\textbf{11.} $g_{\text{月球}} = \dfrac{\num{6.67e-11} \times
\num{7.35e22}}{(\num{1.74e6})^2} = \dfrac{\num{4.90e12}}
{\num{3.03e12}} \approx \qty{1.62}{N/kg}$。

\textbf{12.} 高度 $\propto 1/g$：
$h = 0.50 \times \dfrac{9.81}{1.62} \approx \qty{3.0}{m}$——相当于
慢动作下站立起跳越过篮筐的高度。

\textbf{13.} $g_{\text{水星}} = \dfrac{\num{6.67e-11} \times
\num{3.30e23}}{(\num{2.44e6})^2} = \dfrac{\num{2.20e13}}
{\num{5.95e12}} \approx \qty{3.70}{N/kg}$——几乎恰好等于火星的
\qty{3.73}{N/kg}。水星质量较小，但其半径较小
（即下式中 $R^2$ 较小）抵消了这一点：它是一个密度大得多的星球。

\textbf{14.} 要使\omterm{def:g10:universal-gravitation:fieldstrength}{重力加速度}减半，需满足
$(R + h)^2 = 2 R^2$，因此 $R + h = R\sqrt{2}$，且
$h = (\sqrt{2} - 1) R \approx 0.414 \times \num{6.37e6}
\approx \qty{2.6e6}{m}$——即地面以上约 \qty{2600}{km}。

\textbf{15.} $\dfrac{G M_{\text{地球}}}{d^2} = 1.62$，得
$d = \sqrt{\dfrac{\num{6.67e-11} \times \num{5.97e24}}{1.62}}
= \sqrt{\num{2.46e14}} \approx \qty{1.57e7}{m} \approx 2.5$ 个地球
半径。仅在离地面一个半
地球半径的高度处，地球的引力就不比月球表面更强。

\textbf{16.} $v = \dfrac{2\pi d}{T}
= \dfrac{2\pi \times \num{1.496e11}}{\num{3.156e7}}
\approx \qty{2.98e4}{m/s}$——即每秒三十千米。

\textbf{17.} $\dfrac{v^2}{d} = \dfrac{(\num{2.98e4})^2}{\num{1.496e11}}
\approx \qty{5.9e-3}{N/kg}$：太阳以
约六毫牛的力拉着地球的每\omterm{def:g10:orders-of-magnitude:unit}{千克}物质。

\textbf{18.} $\dfrac{G M_{\text{太阳}}}{d^2} = \dfrac{v^2}{d}$，得
\[
M_{\text{太阳}} = \frac{v^2 d}{G}
= \frac{\num{8.87e8} \times \num{1.496e11}}{\num{6.67e-11}}
\approx \qty{1.99e30}{kg}.
\]

\textbf{19.} $\dfrac{\num{1.99e30}}{\num{5.97e24}} \approx
\num{3.3e5}$：太阳约等于三十三万个地球。

\textbf{20.} 由于同一个 $G$ 支配着每一对质量之间的引力，
在木棚中用两只铅球测定一次它的数值，便能
先在天平上称出脚下的行星，再称出我们所绕行的恒星——
这正是\emph{普适}的价值所在。
\end{solution}
